% ══════════════════════════════════════════════════════════════
%  Chapter 10: Block Acceptance --- "Valid and Not Already Spent"
% ══════════════════════════════════════════════════════════════
\chapter{Block Acceptance}
\label{ch:block-acceptance}

\begin{learnbox}
\begin{itemize}
  \item Explain the whitepaper's Step 5 gate: ``Nodes accept the block only if all transactions in it are valid and not already spent''
  \item Map this safety rule to concrete code boundaries in the implementation
  \item Identify the exact point where ownership and scarcity are enforced
  \item Understand the validate-first, connect-only-if-valid pattern that prevents double spending
\end{itemize}
\end{learnbox}

\lettrine{B}{itcoin's} \index{whitepaper}whitepaper compresses the most important safety rule in the entire system into one sentence:

\begin{quote}
``Nodes accept the \index{block}block only if all \index{transaction}transactions in it are valid and not already spent.''
--- \emph{Bitcoin Whitepaper, Section 5}
\end{quote}

This is the moment where your \index{node}node decides whether incoming data becomes \textbf{\index{durable state}durable state}. Get this wrong and the \index{chain}chain can still ``look like a \index{blockchain}blockchain'' while quietly violating the two rules that make Bitcoin Bitcoin:

\begin{itemize}
  \item \textbf{\index{ownership}Ownership}: only the holder of the right \index{key}key can spend an \index{output}output.
  \item \textbf{\index{scarcity}Scarcity}: each \index{output}output can be spent \textbf{at most once} (no \index{double-spending}double-spends).
\end{itemize}

In the previous \index{chain}chain chapters we learned the pieces:

\begin{itemize}
  \item \textbf{\index{validation}Validity} lives in code like \code{\index{Transaction!verify}Transaction::verify(...)} (\index{digital signature}signatures, structure).
  \item \textbf{\index{spendability}Spendability} lives in the \index{UTXO}UTXO set (what is unspent right now).
  \item \textbf{\index{fork-choice}Fork choice} decides which \index{branch}branch we build on.
\end{itemize}

This capstone is about the \emph{boundary} where those pieces must be composed correctly:

\begin{infobox}[Validate First, Connect Only If Valid]
This chapter isolates the missing piece and shows the exact boundary where it must be implemented in the code.
\end{infobox}

% ──────────────────────────────────────────────────────────────
\section{Scope: Where Step 5 Appears in Our Codebase}
\label{sec:ch10-scope}

In this repository, the ``block became canonical'' boundary is in:

\begin{itemize}
  \item \code{bitcoin/src/store/file\_system\_db\_chain.rs} $\to$ \code{BlockchainFileSystem::add\_block(...)}
\end{itemize}

That's the method that must enforce the Step~5 gate \textbf{before} mutating the canonical tip and the UTXO set.

% ──────────────────────────────────────────────────────────────
\section{The Runtime Story: ``A Block Just Arrived''}
\label{sec:ch10-runtime-story}

Imagine your \index{node}node receives a \code{\index{Package}Package::Block} from a \index{peer}peer. At that moment, you have a choice:

\begin{enumerate}
  \item Treat it as \emph{\index{data}data you can store} (fine: you can persist \index{block}blocks as candidates).
  \item Treat it as \emph{\index{state}state you will build on} (dangerous: this is where Step 5 applies).
\end{enumerate}

The \index{whitepaper}whitepaper's Step 5 is about the second choice:

\begin{quote}
``Accept the \index{block}block'' means ``let this \index{block}block change my \index{node}node's authoritative view of money.''
\end{quote}

That includes:

\begin{itemize}
  \item Moving the \index{canonical tip}canonical tip pointer.
  \item Updating \index{derived state}derived state (\index{UTXO}UTXO set) so ``what is spendable'' matches the new tip.
  \item Potentially \index{reorganization}reorganizing (disconnect old tip \index{block}blocks, connect new \index{branch}branch \index{block}blocks).
\end{itemize}

This is why Step 5 is not a ``nice-to-have validation''; it is the contract that prevents your node from turning invalid history into real balances.

% ──────────────────────────────────────────────────────────────
\section{A Rust Mental Model: The Acceptance Boundary Is an API Boundary}
\label{sec:ch10-mental-model}

As a \index{Rust!developer}Rust developer, a clean way to think about Step 5 is:

\begin{itemize}
  \item \textbf{\index{validation}Validation} is a \emph{\index{pure function}pure function}: it can read \index{chain}chain \index{state}state and \index{UTXO}UTXO \index{state}state, but must not mutate \index{durable state}durable state.
  \item \textbf{\index{connection}Connection} is a \emph{\index{state transition}state transition}: it mutates \index{tip}tip + \index{UTXO}UTXO (and maybe indexes), and must be atomic or \index{rollback}rollback-safe.
\end{itemize}

That gives you a two-phase acceptance interface:

\begin{minted}[fontsize=\footnotesize,linenos=false]{text}
validate(block, view) -> Ok(()) or Err(...)
connect(block, state) -> Ok(()) or Err(...)
\end{minted}

This pattern prevents partial \index{state transition}state transitions from leaving the \index{node}node inconsistent.

% ──────────────────────────────────────────────────────────────
\section{Block Validation Pipeline}
\label{sec:ch10-pipeline}

Here is the diagram you should keep in your head while reading code:

\begin{figure}[htbp]
\centering
\begin{tikzpicture}[
  node distance = 1.0cm,
  stage/.style = {rectangle, rounded corners=4pt, draw=sectioncolor, fill=codebg,
                  text width=2.2cm, align=center, minimum height=0.9cm, font=\sffamily\small},
  check/.style = {font=\tiny\sffamily, color=sectioncolor},
  arrow/.style = {-{Stealth[length=4pt,width=3pt]}, thick, color=brand},
  fail/.style = {-{Stealth[length=4pt,width=3pt]}, thick, dashed, color=codeframe},
]

% Incoming Block
\node[stage] (incoming) at (0, 7) {
  Incoming \\
  Block
};

% Parse Block
\node[stage] (parse) at (0, 5.8) {
  Parse \\
  Header
};

% Check PoW
\node[stage] (pow) at (0, 4.6) {
  Check PoW \\
  (hash < tgt)
};

% Validate Tx Sig
\node[stage] (sig) at (0, 3.4) {
  Validate \\
  Tx Sig
};

% Check UTXO
\node[stage] (utxo) at (0, 2.2) {
  Check UTXO \\
  (not spent?)
};

% Accept / Connect
\node[stage] (accept) at (0, 1.0) {
  Accept / \\
  Connect
};

% Reject nodes
\node[font=\tiny\sffamily\color{codeframe}] (reject1) at (2.5, 5.8) {Reject};
\node[font=\tiny\sffamily\color{codeframe}] (reject2) at (2.5, 4.6) {Reject};
\node[font=\tiny\sffamily\color{codeframe}] (reject3) at (2.5, 3.4) {Reject};
\node[font=\tiny\sffamily\color{codeframe}] (reject4) at (2.5, 2.2) {Reject};

% Main path arrows
\draw[arrow] (incoming) -- (parse);
\draw[arrow] (parse) -- (pow);
\draw[arrow] (pow) -- (sig);
\draw[arrow] (sig) -- (utxo);
\draw[arrow] (utxo) -- (accept);

% Fail arrows
\draw[fail] (parse) -- (reject1);
\draw[fail] (pow) -- (reject2);
\draw[fail] (sig) -- (reject3);
\draw[fail] (utxo) -- (reject4);

% Labels
\node[check] at (-1.5, 6.4) {OK};
\node[check] at (-1.5, 5.2) {OK};
\node[check] at (-1.5, 4.0) {OK};
\node[check] at (-1.5, 2.8) {OK};

\end{tikzpicture}
\caption{Block Validation Pipeline: Parse, PoW, Signatures, and UTXO Checks}
\label{fig:ch10-validation-pipeline}
\end{figure}

If your code discovers invalidity ``mid-update'', you have already violated the contract: you let invalid data partially become state.

% ──────────────────────────────────────────────────────────────
\section{Step 5, Expanded: What ``Valid'' and ``Not Already Spent'' Mean}
\label{sec:ch10-what-valid-means}

The whitepaper is minimal, so we'll be explicit. In a learning implementation like this repo, the smallest useful Step~5 set is:

% ──────────────────────────────────────────────────────────────
\subsection{Block-Level Sanity (Structural)}
\label{sec:ch10-structural}

At minimum:

\begin{itemize}
  \item \textbf{Coinbase rules}: a block should contain \textbf{exactly one coinbase} transaction (simplified rule used throughout the book).
  \item \textbf{Parent linkage}: if we are connecting the block, its parent must be known (or we keep it as an orphan/candidate).
  \item \textbf{Proof-of-work}: if your chain is PoW-secured, the header hash must satisfy the target (some paths in this repo intentionally simplify/todo this).
\end{itemize}

% ──────────────────────────────────────────────────────────────
\subsection{Transaction-Level Validity (Authorization)}
\label{sec:ch10-authorization}

\begin{warnbox}[Critical Gate]
The UTXO check in block validation is the single gate that prevents double spending. Without it, an attacker could include two transactions in the same block that both spend the same output, effectively creating money from nothing.
\end{warnbox}

For each non-coinbase transaction, ``valid'' includes:

\begin{itemize}
  \item \textbf{Signature verification} against the reconstructed digest.
\end{itemize}

In this repo, that primitive already exists:

\begin{itemize}
  \item \code{bitcoin/src/primitives/transaction.rs} $\to$ \code{Transaction::verify(\&BlockchainService)}
\end{itemize}

% ──────────────────────────────────────────────────────────────
\subsection{State-Level Validity (``Not Already Spent'')}
\label{sec:ch10-not-already-spent}

The key phrase is \textbf{``not already spent''}, which is \emph{not} a signature property. It is a statement about the current UTXO view.

This is outpoint-level logic:

\begin{itemize}
  \item \textbf{Outpoint}: $(txid, vout)$
  \item \textbf{Rule}: every input must reference an outpoint that is currently unspent in the UTXO view.
\end{itemize}

Two pitfalls that Step 5 must guard against:

\begin{itemize}
  \item \textbf{Inter-block double spend}: spending an outpoint that was already spent in the current canonical chain.
  \item \textbf{Intra-block double spend}: the same outpoint appears twice as an input across the block's transactions.
\end{itemize}

The intra-block case is easy to miss and easy to fix:

\begin{itemize}
  \item Build a \code{HashSet<(txid, vout)>} while scanning the block.
  \item If you see the same outpoint twice, reject the block immediately.
\end{itemize}

% ──────────────────────────────────────────────────────────────
\section{Mapping Step 5 to Our Repo}
\label{sec:ch10-mapping-to-repo}

The core acceptance entry point is already labeled in the code with a FIXME.

% ──────────────────────────────────────────────────────────────
\subsection{The Acceptance Boundary in Code}
\label{sec:ch10-acceptance-boundary}

\code{bitcoin/src/store/file\_system\_db\_chain.rs} $\to$ \code{BlockchainFileSystem::add\_block(...)}

Inside this method, you'll find a comment like:

\begin{minted}[fontsize=\footnotesize]{rust}
// FIXME: From bitcoin whitepaper, only add block if:
// A) "All transactions in it are valid"
// B) "Not already spent"
\end{minted}

That is the right instinct, and it points to the right place: \textbf{Step 5 belongs in \code{add\_block}}.

The important subtlety is what ``add'' means here:

\begin{itemize}
  \item It's fine to \emph{store the block} as ``data we've heard about.''
  \item It is not fine to \emph{connect the block} (tip + UTXO mutation) without Step~5 validation.
\end{itemize}

% ──────────────────────────────────────────────────────────────
\section{A Concrete Design: Implement ``Validate \texorpdfstring{$\to$}{→} Connect''}
\label{sec:ch10-concrete-design}

This section is intentionally practical. It's not ``Bitcoin Core''; it's a clear shape you can implement in Rust in this codebase.

% ──────────────────────────────────────────────────────────────
\subsection{Phase 1: Validate (Read-Only)}
\label{sec:ch10-phase1-validate}

\textbf{Inputs}:

\begin{itemize}
  \item The block being considered.
  \item A read-only chain API (\code{BlockchainService}) to locate referenced previous transactions (needed for signature verification).
  \item A UTXO view to answer ``is this outpoint unspent?''
\end{itemize}

\textbf{Outputs}:

\begin{itemize}
  \item \code{Ok(())} if Step~5 holds.
  \item A domain error that explains which rule was violated.
\end{itemize}

\textbf{Rust shape} (documentation-level pseudocode):

\begin{listing}[H]
\caption{Step 5 validation pseudocode}
\label{lst:ch10-validate-step5}
\begin{minted}[fontsize=\footnotesize]{rust}
use std::collections::HashSet;

#[derive(Clone, Debug, PartialEq, Eq, Hash)]
pub struct OutPoint {
    pub txid: Vec<u8>,
    pub vout: usize,
}

pub async fn validate_step5(
    block: &Block,
    chain: &BlockchainService,
    utxo: &UTXOSet,
) -> Result<()> {
    // 1) Exactly one coinbase (simplified)
    let txs = block.get_transactions().await?;
    let coinbase_count = txs.iter()
        .filter(|tx| tx.is_coinbase()).count();
    if coinbase_count != 1 {
        return Err(BtcError::InvalidTransaction);
    }

    // 2) No outpoint spent twice in the same block
    let mut seen: HashSet<OutPoint> = HashSet::new();

    for tx in &txs {
        if tx.is_coinbase() {
            continue;
        }

        // 3) Authorization (signatures)
        if !tx.verify(chain).await? {
            return Err(BtcError::InvalidSignature);
        }

        // 4) Spendability (UTXO view)
        for vin in tx.get_vin() {
            let op = OutPoint {
                txid: vin.get_txid().to_vec(),
                vout: vin.get_vout(),
            };

            if !seen.insert(op.clone()) {
                // Intra-block double spend
                return Err(
                    BtcError::InvalidTransactionInput
                );
            }

            // Your UTXO API wants to answer:
            // is (txid, vout) currently unspent?
            utxo.assert_unspent(&op).await?;
        }
    }

    Ok(())
}
\end{minted}
\end{listing}

\begin{notebox}[Key Insight]
Step 5 needs an outpoint-level query (\code{assert\_unspent}), not just ``find UTXOs by owner''.
\end{notebox}

% ──────────────────────────────────────────────────────────────
\subsection{Phase 2: Connect (Mutating)}
\label{sec:ch10-phase2-connect}

Once validation passes, connecting the block is ``just'' the state transition:

\begin{itemize}
  \item Spend the referenced outpoints.
  \item Insert the new outputs.
  \item Advance tip (or apply reorg logic).
\end{itemize}

But there is one critical systems requirement:

\begin{warnbox}[Atomicity Requirement]
Connecting must be atomic or rollback-safe.
\end{warnbox}

If your UTXO update is multiple DB writes and you can fail mid-way, your node can end up in a state that doesn't correspond to any valid chain.

In this repo, the connector is implemented as UTXO update logic (e.g.\ \code{UTXOSet::update(...)}), and the persistence layer is sled. The design goal is:

\begin{itemize}
  \item compute the delta
  \item apply the delta in one sled transaction (or a clearly defined ``write then commit tip'' transaction)
\end{itemize}

% ──────────────────────────────────────────────────────────────
\section{Reorgs: Step 5 Is Not Optional on the ``Apply Blocks'' Path}
\label{sec:ch10-reorgs}

A common misunderstanding is: ``we validated blocks when we first saw them, so reorg apply is safe.''

In a real node, and in this repo's learning node, reorg application still needs the same gate:

\begin{itemize}
  \item When you \textbf{disconnect} blocks: rollback derived state safely.
  \item When you \textbf{connect} blocks on the winning branch: each connected block must satisfy Step 5 against the evolving UTXO view.
\end{itemize}

So any method that effectively does ``connect a sequence of blocks'' must call the same validator you use for normal extension.

In the current codebase, look for reorg helpers like:

\begin{itemize}
  \item \code{reorganize\_chain(...)}
  \item \code{apply\_chain\_from\_ancestor(...)}
\end{itemize}

and treat them as ``batch connect'': Step~5 validation should be enforced per block.

% ──────────────────────────────────────────────────────────────
\section{A Practical Note About UTXO Representation}
\label{sec:ch10-utxo-representation}

Step 5 is an outpoint-level rule. Your UTXO layer must preserve stable outpoint identity:

\begin{itemize}
  \item In Bitcoin, $(txid, vout)$ is stable forever.
  \item If your UTXO storage is \code{txid -> Vec<TXOutput>} and you \emph{remove elements} from the \code{Vec}, you can accidentally shift indices and break the meaning of \code{vout}.
\end{itemize}

That's why many implementations model UTXO state as something like:

\begin{itemize}
  \item \code{outpoint -> TXOutput}
\end{itemize}

or keep a stable vector plus a spent marker rather than physically removing entries.

You don't have to implement a perfect data model for this book, but you do need one property:

\begin{notebox}
\code{vout} must keep its meaning, so ``is $(txid, vout)$ unspent?'' is deterministic.
\end{notebox}

% ──────────────────────────────────────────────────────────────
\section{Conclusion}
\label{sec:ch10-conclusion}

The whitepaper's Step 5 is not a comment, and not ``an error we hope happens during update''.

It is a hard gate:

\begin{itemize}
  \item \textbf{Validate all transactions} (authorization).
  \item \textbf{Reject double spends} (state-level spendability).
  \item \textbf{Only then} mutate the canonical tip and the UTXO set.
\end{itemize}

That's the moment where consensus becomes state.

% ──────────────────────────────────────────────────────────────
\section{Exercises}
\label{sec:ch10-exercises}

\begin{enumerate}
  \item \textbf{Construct a Double-Spend Block} --- Create a block containing two transactions that both attempt to spend the same UTXO. Predict at which validation step the block will be rejected. Trace through the code to verify your prediction.

  \item \textbf{Validation Order Analysis} --- The validate-first pattern checks transactions before connecting the block to the chain. Describe what could go wrong if the order were reversed (connect first, validate second). What specific attack would this enable?
\end{enumerate}

% ──────────────────────────────────────────────────────────────
\begin{chaptersummary}
\item We examined the whitepaper's Step 5 gate --- ``Nodes accept the block only if all transactions in it are valid and not already spent'' --- and mapped it to concrete code boundaries.
\item We traced the validate-first, connect-only-if-valid pattern that prevents double spending at the consensus layer.
\item We identified the exact boundary where ownership and scarcity are enforced: the UTXO check that ensures every input references an unspent output.
\item We connected this capstone validation rule back to the transaction lifecycle, consensus rules, and chain state management covered in previous chapters.
\end{chaptersummary}

In the next chapter, we move from validation to persistence --- how all of this data is stored on disk using the Sled embedded database.
